\chapter{How to use this Flora}

\section*{Arrangement}
Families are arranged alphabetically, and within each family the
genera and species are also arranged alphabetically. The Key to
Families (Chapter~\ref{ch:key-families}) leads to a family; the key at the start
of each family chapter leads to a species.

\section*{Keys}
The keys are dichotomous. Each couplet offers two contrasting leads,
numbered \textbf{1a} and \textbf{1b}, \textbf{2a} and \textbf{2b}, and
so on. Choose the lead that fits the specimen and follow the dotted
line to its end, which gives either the number of the next couplet or
a name. The short example below shows the layout only.

\begin{florakey}
	\lead{1a}{Leaves compound}{2}
	\lead{1b}{Leaves simple}{Key to Families}
	\lead{2a}{Leaves bipinnate; flowers in pendent red heads}{\taxon{Parkia}}
	\lead{2b}{Leaves once pinnate; flowers in panicles}{\taxon{Khaya}}
\end{florakey}

\section*{Names}
Scientific names are printed in italics and are followed by their
author citation in roman type, for example
\sciname{Milicia excelsa}{(Welw.) C.C.Berg}. Author abbreviations
follow the standard forms listed in the Author Abbreviations
appendix. Principal synonyms are given under each accepted name.

\section*{Specimen citations}
Each specimen is cited by collector, collector's number and the
herbarium in which it is kept, for example
\specimen{Alade}{104}{AUH}. Herbarium codes used in this volume are:

\medskip
\begin{tabular}{@{}ll@{}}
	\toprule
	Code & Herbarium \\
	\midrule
	\textsc{auh} & Arami Uplands Herbarium \\
	\textsc{okf} & Okeri Forestry Research Institute Herbarium \\
	\textsc{svn} & Savanna Botanical Station Herbarium \\
	\bottomrule
\end{tabular}

\section*{Abbreviations}
\begin{tabular}{@{}ll@{}}
	alt.  & altitude \\
	c.    & about (\textit{circa}) \\
	fl.   & flowering \\
	fr.   & fruiting \\
	s.n.  & without number (\textit{sine numero}) \\
	st.   & sterile (neither flowers nor fruits) \\
	syn.  & synonym \\
\end{tabular}